\documentclass[10pt,a4paper]{amsart}
\usepackage[margin=19mm]{geometry}
\usepackage{amsmath,amssymb,mathtools}
\usepackage[scr=rsfs,cal=euler]{mathalfa}
\usepackage{xfrac}
\usepackage[unicode,hidelinks,bookmarks=false]{hyperref}
\newcommand{\ie}{i.e.\ }
\newcommand{\cf}{cf.\ }
\newcommand{\resp}{respectively\ }
\newcommand{\Subsection}{\subsection}
\providecommand{\nobreakdash}{\nobreak\hbox{-}}
\setlength{\emergencystretch}{2em}
\multlinegap0pt
\usepackage{amssymb}
\usepackage[shortlabels]{enumitem}
\usepackage{mathtools}
\newtheorem{lemma}{Lemma}
\title{Parabolic gap theorems for the Yang-Mills energy}
\author{Anuk Dayaprema, Alex Waldron}
\date{Source: arXiv:2412.21050v2 (CC BY 4.0)}
\begin{document}
\maketitle

\noindent\emph{Source and licence.} Parabolic gap theorems for the Yang-Mills energy. Version arXiv:2412.21050v2 (CC BY 4.0), distributed by arXiv under the Creative Commons Attribution 4.0 International licence, http://creativecommons.org/licenses/by/4.0/, as recorded in the arXiv licence metadata for this exact version and shown on the abstract page https://arxiv.org/abs/2412.21050v2. Full licence notice: \texttt{licenses/arxiv-2412.21050-CC-BY-4.0.txt}.\par\smallskip

\noindent\textit{Editorial scope.} Counted unit: the article's lemma lemma:comparablesup (source0.tex lines 1134--1138). The article's complete original proof of it is delivered with it (source0.tex lines 1139--1155, one \texttt{proof} environment of 17 lines). No mathematics is written, repaired or re-derived: the statement and the proof are the article's own lines, in the article's own notation, and no retained line is substituted or re-flowed.  No further source-local context is needed: every cross-reference of every retained line resolves inside the two delivered spans.  Nothing was omitted from any retained span, so the package carries no omission record.  No retained line contains a citation, so the wrapper carries no bibliography and no bibliography entry.  No retained line contains a figure environment, an image-inclusion command or drawing code, so no image is delivered and the package declares no asset dependency.  Editorial formatting only, in addition: an AMS \texttt{amsart} preamble that restates the article's own declarations and macros used by the retained lines, namely the environments \texttt{lemma} on separate counters, together with the standard packages those lines need.  The retained environments print in this document as Lemma 1 in order of appearance, whereas the article prints the same items with the numbering of its own full document.  The complete native source of the article as distributed in the arXiv e-print for this exact version is bundled unchanged as \texttt{sources/170925/source0.tex}.
\begin{lemma}\label{lemma:comparablesup}
    Let \(\lambda > 0\) and let \(f\) be a non-negative continuous function on \(P_r(x_0, t_0)\) such that \(f(x_0, t_0) = \lambda^{-2} > 16r^{-2}\). Then there exists a spacetime point \((\tilde{x}, \tilde{t})\) such that, denoting \(\mu^{-2} := f(\tilde{x}, \tilde{t})\), we have \(\mu \leq \lambda\), \((\tilde{x}, \tilde{t}) \in P_{2\lambda}(x_0, t_0)\), and     \begin{align} \label{comparablesupbound}
\sup_{P_\mu(\tilde{x}, \tilde{t})} f \leq 4\mu^{-2}.
    \end{align}
\end{lemma}

\begin{proof}
    Suppose for %a 
    contradiction that such a spacetime point does not exist. Then in particular the desired conclusion does not hold for \((x_0, t_0)\), so we have
    \begin{align*}
        \sup_{P_\lambda(x_0, t_0)} f > 4\lambda^{-2}.
    \end{align*} Assume by induction that we have found \(n+1\) spacetime points \((x_i, t_i)\), \(i = 0, \ldots, n\) satisfying
    the following conditions:
    \begin{align*}
       (\ast) \begin{cases}
            (x_{i+1}, t_{i+1}) \in \overline{P_{\mu_i}(x_i, t_i)} \ \ \ \text{ where } \mu_i^{-2} := f(x_i, t_i), \\
            f(x_{i + 1}, t_{i+1}) = \sup_{P_{\mu_i}(x_i, t_i)} f > 4\mu_i^{-2}.
        \end{cases}
    \end{align*} Note that \(\mu_{i+1} < 2^{-1}\mu_i\), and hence \(\mu_n < 2^{-n}\lambda\). We let \(d\) denote the parabolic distance \(d((p, t), (q, s)) = \max\{d_g(p, q), \sqrt{|t-s|}\}\), for which parabolic cylinders are metric balls. Observe that 
    \begin{align*}
        d((x_0, t_0), (x_n, t_n)) \leq \sum_{i = 0}^n d((x_{i+1}, t_{i+1}), (x_i, t_i)) \leq \sum_{i = 0}^n \mu_i < \lambda \sum_{i = 0}^n 2^{-i} < 2\lambda.
    \end{align*} Thus for \(0 \leq i \leq n-1\), \((x_i, t_i)\) is in \(P_{2\lambda}(x_0, t_0)\). Additionally, \(P_{\mu_i}(x_i, t_i)\) is compactly contained in \(P_r(x_0, t_0)\) since \(\lambda < r/4\). Our assumption that there is no spacetime point satisfying the conclusion of the lemma implies the existence of a spacetime point \((x_{n+1}, t_{n+1})\) satisfying \((\ast)\). Therefore, by induction we obtain an infinite sequence of spacetime points \(\{(x_i, t_i)\}\) satisfying \((\ast)\). Since this sequence is contained in the compact set \(\overline{P_{2\lambda}(x_0, t_0)}\) and since \(f\) is continuous, and yet \(\mu_i^{-2}\) tends to infinity, we reach a contradiction, as desired.
\end{proof}

\end{document}
